
\documentclass{conm-p-l}
%%%%%%%%%%%%%%%%%%%%%%%%%%%%%%%%%%%%%%%%%%%%%%%%%%%%%%%%%%%%%%%%%%%%%%%%%%%%%%%%%%%%%%%%%%%%%%%%%%%%%%%%%%%%%%%%%%%%%%%%%%%%%%%%%%%%%%%%%%%%%%%%%%%%%%%%%%%%%%%%%%%%%%%%%%%%%%%%%%%%%%%%%%%%%%%%%%%%%%%%%%%%%%%%%%%%%%%%%%%%%%%%%%%%%%%%%%%%%%%%%%%%%%%%%%%%
\usepackage{amsfonts}

\setcounter{MaxMatrixCols}{10}
%TCIDATA{OutputFilter=LATEX.DLL}
%TCIDATA{Version=5.50.0.2960}
%TCIDATA{<META NAME="SaveForMode" CONTENT="1">}
%TCIDATA{BibliographyScheme=Manual}
%TCIDATA{Created=Thursday, February 25, 2021 13:05:24}
%TCIDATA{LastRevised=Friday, February 26, 2021 08:45:04}
%TCIDATA{<META NAME="GraphicsSave" CONTENT="32">}
%TCIDATA{<META NAME="DocumentShell" CONTENT="Author Packages for AMS\Contemporary Mathematics (Proceedings)">}
%TCIDATA{CSTFile=conm-p-l.cst}

\newtheorem{theorem}{Theorem}[section]
\newtheorem{lemma}[theorem]{Lemma}
\theoremstyle{definition}
\newtheorem{definition}[theorem]{Definition}
\newtheorem{example}[theorem]{Example}
\newtheorem{xca}[theorem]{Exercise}
\theoremstyle{remark}
\newtheorem{remark}[theorem]{Remark}
\numberwithin{equation}{section}
\theoremstyle{plain}
\newtheorem{acknowledgement}{Acknowledgement}
\newtheorem{algorithm}{Algorithm}
\newtheorem{axiom}{Axiom}
\newtheorem{case}{Case}
\newtheorem{claim}{Claim}
\newtheorem{conclusion}{Conclusion}
\newtheorem{condition}{Condition}
\newtheorem{conjecture}{Conjecture}
\newtheorem{corollary}{Corollary}
\newtheorem{criterion}{Criterion}
\newtheorem{exercise}{Exercise}
\newtheorem{notation}{Notation}
\newtheorem{problem}{Problem}
\newtheorem{proposition}{Proposition}
\newtheorem{solution}{Solution}
\newtheorem{summary}{Summary}
\copyrightinfo{2001}{enter name of copyright holder}

\input{tcilatex}

\begin{document}
\title[Short Title]{Contemporary Mathematics (Proceedings)}
\author{Author One}
\address[A. One and A. Two]{Author OneTwo address line 1\\
Author OneTwo address line 2}
\email[A. One]{aone@aoneinst.edu}
\urladdr{http://www.authorone.oneuniv.edu}
\thanks{Thanks for Author One.}
\author{Author Two}
\curraddr[A. Two]{Author Two current address line 1\\
Author Two current address line 2}
\email[A.~Two]{atwo@atwoinst.edu}
\urladdr{http://www.authortwo.twouniv.edu}
\thanks{Thanks for Author Two.}
\author{Author Three}
\address[A. Three]{Author Three address line 1\\
Author Three address line 2}
\urladdr{http://www.authorthree.threeuniv.edu}
\thanks{This paper is in final form and no version of it will be submitted
for publication elsewhere.}
\subjclass[2000]{Primary 05C38, 15A15; Secondary 05A15, 15A18}
\date{July 27, 2001}
\keywords{Keyword one, keyword two, keyword three}
\dedicatory{Dedicated to the memory of S. Bach.}

\begin{abstract}
Replace this text with your own abstract.
\end{abstract}

\maketitle

\section{Lattices as partially ordered sets}

If $\left\{ L,\leq \right\} $ is a partially ordered set \emph{(poset)}, and 
$S\subseteq L$ is an arbitrary subset, then an element $u\in L$ is said to
be an \emph{upper bound} of $S$ if 
\begin{equation}
s\leq u\text{ for each }s\in S.
\end{equation}
\emph{A set may have many upper bounds, or none at all}. An upper bound $u$
of $S$ is said to be its \emph{least upper bound}, or \emph{join}, or \emph{%
supremum}, if $u\leq x$ for each upper bound $x$ of $S$. \emph{A set need
not have a least upper bound}, but it cannot have more than one. Dually, $%
l\in L$ is said to be a lower bound of $S$ if $l\leq s$ for each $s\in S$. A
lower bound $l$ of $S$ is said to be its \emph{greatest lower bound}, or 
\emph{meet}, or \emph{infimum}, if$x$ $\leq l$ for each lower bound $x$ of $S
$. \emph{A set may have many lower bounds, or none at all}, but can have at
most one greatest lower bound.

A partially ordered set $(L,\leq )$ is called a \emph{join-semilattic}e if
each two-element subset $\{a,b\}\subseteq L$ has a join (i.e. least upper
bound), and is called a\emph{\ meet-semilattice} if each two-element subset
has a meet (i.e. greatest lower bound), denoted by $a\vee b$ and $a\wedge b$
respectively. $\left\{ L,\leq \right\} $ is called a \emph{lattice} if it is 
\emph{both a join- and a meet-semilattice.} This definition makes $\vee $
and $\wedge $ binary operations. Both operations are monotone with respect
to the given order:%
\begin{equation}
a1\leq a2\quad \&\quad b1\leq b2\Rightarrow a1\vee b1\leq a2\vee b2\quad
\&\quad a1\wedge b1\leq a2\wedge b2.
\end{equation}
It follows by an induction argument that every non-empty finite subset of a
lattice has a \emph{least upper bound} and a \emph{greatest lower bound}. 

A \emph{bounded lattice} is a lattice that \emph{additionally} has a \emph{%
greatest element (also called maximum, or top) element}, and denoted by $1$,
or by $\U{22a4} $  and a \emph{least element (also called minimum, or
bottom),} denoted by $0$ or by $\perp $, which satisfy%
\begin{equation}
0\leq x\leq 1\forall x\in L.
\end{equation}

Every lattice can be embedded into a bounded lattice by adding an artificial
greatest and least element, and every non-empty finite lattice is bounded,
by taking the join (respectively, meet) of all elements, denoted by%
\begin{equation}
\U{22c1} L=a_{1}\vee \cdots \vee a_{n}
\end{equation}%
and 
\begin{equation}
\U{22c0} L=a_{1}\wedge \cdots \wedge a_{n}.
\end{equation}
where%
\begin{equation}
L=\{a_{1},\ldots \ ,a_{n}\}.
\end{equation}

A partially ordered set is a bounded lattice if and only if every finite set
of elements (including the empty set) has a join and a meet. For every
element $x$ of a poset it is trivially true that%
\begin{eqnarray}
\forall a &\in &\emptyset :x\leq a \\
\forall a &\in &\emptyset :x\geq a
\end{eqnarray}
therefore every element of a poset is both an upper bound and a lower bound
of the empty set. This implies that the join of an empty set is the least
element 
\begin{equation}
\U{22c1} \emptyset =0
\end{equation}
and the meet of the empty set is the greatest element 
\begin{equation}
\U{22c0} \emptyset =1
\end{equation}%
This is consistent with the associativity and commutativity of meet and
join: the join of a union of finite sets is equal to the join of the joins
of the sets, and dually, the meet of a union of finite sets is equal to the
meet of the meets of the sets, i.e., for finite subsets $A$ and $B$ of a
poset $L$,

\begin{equation}
\U{22c1} (A\cup B)=(\U{22c1} A)\vee (\U{22c1} B)
\end{equation}

and

\begin{equation}
\U{22c0} (A\cup B)=(\U{22c0} A)\wedge (\U{22c0} B)
\end{equation}

hold. Taking $B$ to be the empty set,%
\begin{equation}
\U{22c1} (A\cup \emptyset )=(\U{22c1} A)\vee (\U{22c1} \emptyset )=(\U{22c1}
A)\vee 0=\U{22c1} A
\end{equation}

and%
\begin{equation}
\U{22c0} (A\cup \emptyset )=(\U{22c0} A)\wedge (\U{22c0} \emptyset )=(%
\U{22c0} A)\wedge 1=\U{22c0} A,
\end{equation}

which is consistent with the fact that 
\begin{equation}
A\cup \emptyset =A.
\end{equation}

A lattice element $y$ is said to cover another element $x$, if $y>x$, but
there does not exist a $z$ such that $y>z>x$. Here, $y>x$ means $x\leq y$
and $x\neq y.$

A lattice $\left\{ L,\leq \right\} $ is called \emph{graded}, sometimes 
\emph{ranked}, if it can be equipped with a rank function $r$ from $L$ to $%
%TCIMACRO{\U{2115} }%
%BeginExpansion
\mathbb{N}
%EndExpansion
$, sometimes to $%
%TCIMACRO{\U{2124} }%
%BeginExpansion
\mathbb{Z}
%EndExpansion
$, compatible with the ordering (so $r(x)<r(y)$ whenever $x<y$) such that
whenever $y$ covers $x$, then $r(y)=r(x)+1$. The value of the rank function
for a lattice element is called its \emph{rank}.

Given a subset of a lattice, $H\subseteq L$, meet and join restrict to
partial functions -- they are undefined if their value is not in the subset $%
H$. The resulting structure on $H$is called a \emph{partial lattice}. In
addition to this extrinsic definition as a subset of some other algebraic
structure (a lattice), a partial lattice can also be\emph{\ intrinsically
defined} as a set with two partial binary operations satisfying certain
axioms.

\section{Lattices as algebraic structures}

\subsection{General lattice}

An algebraic structure $\left\{ L,\vee ,\wedge \right\} $, consisting of a
set $L$  and two binary, commutative and associative operations $\vee $  ,
and $\wedge $ $L$ is a \emph{lattice} if the following axiomatic identities
hold for all elements $a,b\in L$, sometimes called absorption laws.

\begin{enumerate}
\item $a\vee (a\wedge b)=a$ 

\item $a\wedge (a\vee b)=a$ 

The following two identities are also usually regarded as axioms, even
though they follow from the two absorption laws taken together. Those are
called idempotent laws.

\item $a\vee a=a$

\item $a\wedge a=a$
\end{enumerate}

These axioms assert that both $\left\{ L,\vee \right\} $ and $\left\{
(L,\wedge )\right\} $ are semilattices. The absorption laws, the only axioms
above in which both meet and join appear, distinguish a lattice from an
arbitrary pair of semilattice structures and assure that the two
semilattices interact appropriately. In particular, each semilattice is the
dual of the other.

\subsubsection{Bounded lattice}

A bounded lattice is an algebraic structure of the form $\left\{ (L,\vee
,\wedge ,0,1)\right\} $ such that $\left\{ (L,\vee ,\wedge )\right\} $is a
lattice, $0$ (the lattice's bottom) is the identity element for the join
operation  , and $1$ (the lattice's top) is the identity element for the
meet operation $\wedge $

\begin{enumerate}
\item $a\vee 0=a$

\item $a\wedge 1=a$
\end{enumerate}

\subsubsection{Connection between the two definitions}

An order-theoretic lattice gives rise to the two binary operations $\vee $
and $\wedge $. Since the commutative, associative and absorption laws can
easily be verified for these operations, they make $\left\{ L,\vee ,\wedge
\right\} $ into a lattice in the algebraic sense.

The converse is also true. Given an algebraically defined lattice$\left\{
L,\vee ,\wedge \right\} $, one can define a partial order$\leq $ on $L$ by
setting

\begin{equation}
a\leq bifa=a\wedge b,
\end{equation}
or%
\begin{equation}
a\leq bifb=a\vee b
\end{equation}%
for all elements $\forall a,b\in L$. The laws of absorption ensure that both
definitions are equivalent:

\begin{equation}
a=a\wedge b\Rightarrow b=b\vee (b\wedge a)=(a\wedge b)\vee b=a\vee b
\end{equation}%
and dually for the other direction. One can now check that the relation $%
\leq $ introduced in this way defines a partial ordering within which binary
meets and joins are given through the original operations $\vee $ and $%
\wedge $.

Since the two definitions of a lattice are equivalent, one may freely invoke
aspects of either definition in any way that suits the purpose at hand. 

\end{document}
